% ===== 原文 pp.184–198（PDF pp.191–205） =====
% [原文 p.184]
\begin{proposition}{7.33}
对 $A_\infty$ 代数 $(A,\mathbf m)$、$(A',\mathbf m')$、$(A'',\mathbf m'')$ 与 $A_\infty$ 写像 $\mathbf f\colon(A,\mathbf m)\to(A',\mathbf m')$、$\mathbf g\colon(A',\mathbf m')\to(A'',\mathbf m'')$，有
\[
  \mathcal M(\mathbf g\circ\mathbf f)=\mathcal M(\mathbf g)\circ\mathcal M(\mathbf f).
\]
\end{proposition}

\begin{theorem}{7.34}
设 $\mathbf f,\mathbf g\colon(A,\mathbf m)\to(A',\mathbf m')$ 是 $A_\infty$ 代数 $(A,\mathbf m)$、$(A',\mathbf m')$ 之间的 $A_\infty$ 写像。若 $\mathbf f$ 与 $\mathbf g$ 是 $A_\infty$ 同伦的，则 $\mathcal M(\mathbf f)\cong\mathcal M(\mathbf g)$。
\end{theorem}

\begin{proof}
只需对每个 $a\in\mathcal M(A,\mathbf m)$ 证明 $f(a)\overset{\mathrm{g.e.}}{\sim}g(a)$ 即可。只要考虑存在 $A_\infty$ 写像 $\tilde{\mathbf f}\colon(\tilde A,\tilde{\mathbf m})\to(\tilde A',\tilde{\mathbf m}')$ 使 $\tilde{\mathbf f}|_{t=0}=\mathbf f$、$\tilde{\mathbf f}|_{t=1}=\mathbf g$ 的情形就够了。由补题 7.31，$\tilde f(a)\in\mathrm{MC}(\tilde A[1],\tilde{\mathbf m}')$，于是可知 $f(a)\overset{\mathrm{g.e.}}{\sim}g(a)$。\hfill$\square$
\end{proof}

\begin{corollary}{7.35}
设 $(A,\mathbf m)$、$(A',\mathbf m')$ 为 $A_\infty$ 代数。若存在 $A_\infty$ 拟同构写像 $\mathbf f\colon(A,\mathbf m)\to(A',\mathbf m')$，则集合 $\mathcal M(A,\mathbf m)$ 与 $\mathcal M(A',\mathbf m')$ 之间存在双射。
\end{corollary}

\begin{proof}
由定理 7.26，$\mathbf f$ 是 $A_\infty$ 同伦等价写像，故存在 $A_\infty$ 写像 $\mathbf g\colon(A',\mathbf m')\to(A,\mathbf m)$ 使得
\[
  \mathbf g\circ\mathbf f\overset{\mathrm{h.e.}}{\sim}\id_{(A,\mathbf m)},\qquad \mathbf f\circ\mathbf g\overset{\mathrm{h.e.}}{\sim}\id_{(A',\mathbf m')}.
\]
由补题 7.33 与定理 7.34 得
\[
  \mathcal M(\mathbf g)\circ\mathcal M(\mathbf f)=\id_{\mathcal M(A,\mathbf m)},\qquad \mathcal M(\mathbf f)\circ\mathcal M(\mathbf g)=\id_{\mathcal M(A',\mathbf m')},
\]
于是可知 $\mathcal M(\mathbf f)$、$\mathcal M(\mathbf g)$ 是双射。\hfill$\square$
\end{proof}

\section{$A_\infty$ 范畴}

$A_\infty$ 范畴，是把 $A_\infty$ 代数的元素提升为范畴的态射而成的东西。

\begin{definition}{7.36}
$A_\infty$ 范畴（$A_\infty$-category）$\mathcal C$ 是指：给定对象的集合（或对象的类）$\Ob(\mathcal C)=\{X,Y,\dots\}$，对任意两个对象 $X,Y\in\mathcal C$ 给定 $\Z$ 分次向量空间 $\mathcal C(X,Y)$，并进一步给定次数为 $(2-n)$ 的多重线性映射
\[
  m_{n}\colon\mathcal C(X_{1},X_{2})\otimes\cdots\otimes\mathcal C(X_{n},X_{n+1})\longrightarrow\mathcal C(X_{1},X_{n+1}),\qquad X_{1},\dots,X_{n+1}\in\Ob(\mathcal C)
\]
所成的族 $\mathbf m:=\{m_{n}\}_{n\ge 1}$，且它们确定 $A_\infty$ 结构者。
\end{definition}

在 $A_\infty$ 范畴 $\mathcal C$ 中，把其 $m_{2}$ 视为态射的合成时，$\mathcal C$ 并不成为通常意义下的范畴。实际上当 $m_{3}\ne 0$ 时，一般态射的合成不满足结合律。而且，对每个对象 $X\in\mathcal C$，也不假定恒等元 $1_{X}$ 的存在。

\begin{definition}{7.37}
$A_\infty$ 范畴 $\mathcal C$ 若满足：对每个对象 $X\in\mathcal C$ 存在次数为零的态射 $1_{X}\in\mathcal C^{0}(X,X)$，且对任意态射 $\phi_{12},\phi_{23},\dots$（$\phi_{i(i+1)}\in\mathcal C(X_{i},X_{i+1})$）都有
\[
  m_{2}(1_{X_{1}},\phi_{12})=\phi_{12}=m_{2}(\phi_{12},1_{X_{2}}),\qquad
  m_{n}(\phi_{12},\dots,1_{X_{i}},\dots,\phi_{n(n+1)})=0\quad(n\ge 3),
\]
则称 $\mathcal C$ 具有恒等态射（unital）。
\end{definition}

% [原文 p.185]
$A_\infty$ 范畴 $\mathcal C$ 中高阶积为平凡（$m_{3}=m_{4}=\cdots=0$）者，称为（$K$ 上的）DG 范畴（DG-category）。定义 5.49 中的 DG 范畴，只要把记法按
\[
  \mathcal C(Y,X)=\Hom_{\mathcal C}(X,Y),\qquad m_{2}(\phi,\phi')=m(\phi,\phi')
\]
对齐，就与这里一致。不过，定义 5.49 把恒等态射的存在列入了定义，而这里不作此假定。以后凡具有恒等态射者，都会明确说明。

与 $A_\infty$ 代数的情形一样（参见前节），对任意的 $X,Y\in\mathcal C$，在上同调 $H\bigl(\mathcal C(X,Y),m_{1}\bigr)$ 上，由 $m_{2}$ 诱导出（结合的）合成。

\begin{definition}{7.38}
（$A_\infty$ 范畴的上同调）对 $A_\infty$ 范畴 $\mathcal C$，如下定义（不一定具有恒等态射的）$\Z$ 分次范畴 $H(\mathcal C)$：
\begin{itemize}[leftmargin=1.6em,itemsep=1pt,topsep=1pt]
  \item $\Ob(H(\mathcal C)):=\Ob(\mathcal C)$，
  \item 对每个对象 $X,Y\in\mathcal C$ 与 $r\in\Z$，次数为 $r$ 的态射空间定义为
  \[
    H^{r}\bigl(\mathcal C(X,Y)\bigr):=H^{r}\bigl(\mathcal C(X,Y),m_{1}\bigr).
  \]
\end{itemize}
称 $H(\mathcal C)$ 为 $A_\infty$ 范畴 $\mathcal C$ 的上同调。特别地，由 $H(\mathcal C)$ 只取态射空间为次数零者所得的范畴，称为 $\mathcal C$ 的零次上同调，记作 $H^{0}(\mathcal C)$。
\end{definition}

\begin{definition}{7.39}
$A_\infty$ 范畴 $\mathcal C$ 的悬垂（suspension）$s\mathcal C$ 是指：对象保持不变，即 $\Ob(\mathcal C)=\Ob(s\mathcal C)$，而把态射的次数整体加 $1$ 者：
\[
  (s\mathcal C)^{i}(X,Y):=\mathcal C^{i+1}(X,Y),\qquad X,Y\in\Ob(\mathcal C)=\Ob(s\mathcal C).
\]
\end{definition}

由此，$s\mathcal C$ 上 $A_\infty$ 结构 $m_{k}$（$k=1,2,\dots$）的次数不论 $k$ 都为 $1$。

其次，给定 $A_\infty$ 范畴 $\mathcal C$ 时，对任意的有限集合 $S\subset\Ob(\mathcal C)$，$\Z$ 分次向量空间 $\bigoplus_{X,Y\in S}\mathcal C(X,Y)$ 与 $\{m_{n}\}_{n\ge 1}$ 一起构成 $A_\infty$ 代数。

作为两个 $A_\infty$ 代数之间 $A_\infty$ 写像的类似物，可定义 $A_\infty$ 范畴之间的函子。

\begin{definition}{7.40}
两个 $A_\infty$ 范畴 $\mathcal C,\mathcal C'$ 之间的 $A_\infty$ 函子 $\mathbf f:=\{f;f_{1},f_{2},\dots\}\colon\mathcal C\to\mathcal C'$ 是指：写像 $f\colon\Ob(\mathcal C)\to\Ob(\mathcal C')$，与对任意 $X_{1},\dots,X_{k+1}\in\Ob(\mathcal C)$ 的次数 $(1-k)$ 的多重线性写像
\[
  f_{k}\colon\mathcal C(X_{1},X_{2})\otimes\cdots\otimes\mathcal C(X_{k},X_{k+1})\longrightarrow\mathcal C'\bigl(f(X_{1}),f(X_{k+1})\bigr)
\]
所成的族，且对任意有限集合 $S\subset\Ob(\mathcal C)$，$\{f_{k}\}_{k\ge 1}$ 构成从 $A_\infty$ 代数 $\bigoplus_{X,Y\in S}\mathcal C(X,Y)$ 到 $A_\infty$ 代数 $\bigoplus_{X',Y'\in f(S)}\mathcal C'(X',Y')$ 的 $A_\infty$ 写像。
\end{definition}

\begin{definition}{7.41}
$A_\infty$ 函子 $\mathbf f\colon\mathcal C\to\mathcal C'$ 若对所有的对象 $X\in\mathcal C$ 满足
\[
  f_{1}(\id_{X})=\id_{f(X)},\qquad f_{k}(\dots,\id_{X},\dots)=0,\quad k\ge 2,
\]
则称其保持恒等态射。
\end{definition}

% [原文 p.186]
$m_{1}$ 平凡的 $A_\infty$ 范畴，称为极小 $A_\infty$ 范畴（minimal $A_\infty$-category）。

\begin{definition}{7.42}
对 $A_\infty$ 范畴 $\mathcal C,\mathcal C'$，$A_\infty$ 函子 $\mathbf f\colon\mathcal C\to\mathcal C'$ 若满足 $f\colon\Ob(\mathcal C)\to\Ob(\mathcal C')$ 是双射，且对所有有限集合 $S\subset\Ob(\mathcal C)$，$\mathbf f$ 所定义的 $A_\infty$ 写像是 $A_\infty$ 拟同构写像，则称为 $A_\infty$ 拟同构函子（$A_\infty$-quasi-isomorphism）。
\end{definition}

与 $A_\infty$ 代数有关的极小模型定理（定理 7.17）、分解定理（定理 7.21）等同样的定理，对任意的（具有恒等态射的）$A_\infty$ 范畴 $\mathcal C$ 也成立。$A_\infty$ 范畴 $\mathcal C$ 的极小模型 $\mathcal C'$ 是指与 $\mathcal C$ 为 $A_\infty$ 拟同构的极小 $A_\infty$ 范畴。$\mathcal C$ 的上同调 $H(\mathcal C)$（定义 7.38）与在极小模型 $\mathcal C'$ 中忘掉其高阶 $A_\infty$ 积 $m'_{3},m'_{4},\dots$ 所得的范畴同构。

\begin{definition}{7.43}
若 $A_\infty$ 范畴 $\mathcal C$ 与 $H(\mathcal C)$ 为 $A_\infty$ 拟同构，则称 $\mathcal C$ 是形式的（formal）。
\end{definition}

与 $A_\infty$ 代数的情形一样，也可以定义 $A_\infty$ 范畴之间的 $A_\infty$ 同伦等价函子；在特征零的体 $K$ 上的 $A_\infty$ 范畴的情形，$A_\infty$ 拟同构函子与 $A_\infty$ 同伦等价函子是同一回事。

\section{三角 $A_\infty$ 范畴}

以下在本章中，所谓 $A_\infty$ 范畴，都指某个体 $K$ 上的、具有恒等元的 $A_\infty$ 范畴。

\subsection{由 $A_\infty$ 范畴构造三角范畴的方法}

作为 5.5.2 节中 DG 范畴情形的推广，对任意 $A_\infty$ 范畴 $\mathcal C$，有构造三角范畴的方法（Kontsevich [37]）。也就是说，对 $A_\infty$ 范畴 $\mathcal C$，以 $\mathcal C$ 的单侧扭转复形所构成的 $A_\infty$ 范畴 $Tw(\mathcal C)$ 的零次上同调（定义 7.38），可构成三角范畴 $Tr(\mathcal C)$。
\[
\begin{tikzcd}[row sep=1.2em, column sep=1.2em]
  \{\ A_\infty\ \text{范畴}\ \mathcal C\ \} \arrow[r] & \{\ A_\infty\ \text{范畴}\ Tw(\mathcal C)\ \} \arrow[d,"H^{0}"]\\
  & \{\ \text{三角范畴}\ Tr(\mathcal C)\ \}
\end{tikzcd}
\]
这个 $A_\infty$ 范畴 $Tw(\mathcal C)$ 与三角范畴 $Tr(\mathcal C)$ 按以下步骤构成。

\begin{definition}{7.44}
（$A_\infty$ 范畴 $\tilde{\mathcal C}$）对 $A_\infty$ 范畴 $\mathcal C$，定义 $A_\infty$ 范畴 $\tilde{\mathcal C}$：以 $\mathcal C$ 的对象的次数平移副本及其直和（以及零对象）为对象：
\[
  \Ob(\tilde{\mathcal C})=\Bigl\{\textstyle\bigoplus_{i=1}^{l}\tilde X_{i}\ \Bigm|\ \tilde X_{i}=X_{i}[r_{i}],\ X_{i}\in\mathcal C,\ r_{i}\in\Z\Bigr\}.
\]
这里，态射空间对 $X,Y\in\mathcal C$ 把它们的次数分别减少 $n,m$ 后定义为
\[
  \tilde{\mathcal C}^{r}\bigl(X[n],Y[m]\bigr):=\mathcal C^{r+(m-n)}(X,Y).\]
如此定义，并对直和自然延拓。由 $\mathcal C$ 的 $A_\infty$ 结构 $\mathbf m$ 自然确定的 $\tilde{\mathcal C}$ 的 $A_\infty$ 结构记作 $\tilde m_{1},\tilde m_{2},\dots$。
\end{definition}

% [原文 p.187]
\begin{definition}{7.45}
对对象 $\mathcal X\in\tilde{\mathcal C}$，若能定义次数为零的态射 $\Phi\in(s\tilde{\mathcal C})^{0}(\mathcal X,\mathcal X)$ 使其满足 $s\tilde{\mathcal C}(\mathcal X,\mathcal X)$ 的 $A_\infty$ Maurer--Cartan 方程
\[
  \tilde m_{1}(\Phi)+\tilde m_{2}(\Phi,\Phi)+\tilde m_{3}(\Phi,\Phi,\Phi)+\cdots=0,
\]
则称 $(\mathcal X,\Phi)$ 为 $\mathcal C$ 的扭转复形（twisted complex）。$\mathcal C$ 的单侧扭转复形（one-sided twisted complex）$(\mathcal X,\Phi)$ 是指扭转复形中，对 $\mathcal X=\bigoplus_{i=1}^{l}\tilde X_{i}$ 把 $\Phi=\{\phi_{ij}\}_{i,j=1,\dots,l}$（$\phi_{ij}\in(s\tilde{\mathcal C})^{0}(\tilde X_{i},\tilde X_{j})$）表示时，若 $i\ge j$ 则 $\phi_{ij}=0$ 者。
\end{definition}

\begin{remark}{7.46}
当 $\mathcal C$ 的 $A_\infty$ 积 $m_{n}$ 对充分大的 $n$ 并不平凡时，上述 $A_\infty$ Maurer--Cartan 方程含有无穷和，故扭转复形 $(\mathcal X,\Phi)$ 未必总能定义。然而对于单侧扭转复形 $(\mathcal X,\Phi)$，只要 $\Phi$ 是对角成分为零的上三角矩阵，$m_{n}(\Phi,\dots,\Phi)$ 对大于矩阵尺寸的 $n$ 就自动为零，因此 $A_\infty$ Maurer--Cartan 方程总是无矛盾地有定义。
\end{remark}

% [Tutor Note: 注意 7.46 的「上三角」条件是本书的技术核心：正因为 $\Phi$ 严格上三角，$m_n(\Phi,\dots,\Phi)$ 中任何一项都含足够多的 $\Phi$ 因子而必然落在 $0$ 上，无穷和才截断。若允许 $\Phi$ 有对角块，就必须处理无穷和即「扭转到收敛」的问题，这正是 $\mathrm{Tw}$ 与「形式扭转复形」的区别。]

\begin{definition}{7.47}
（$A_\infty$ 范畴 $Tw(\mathcal C)$）给定 $A_\infty$ 范畴 $\mathcal C$，如下定义由 $\mathcal C$ 的单侧扭转复形所构成的 $A_\infty$ 范畴 $Tw(\mathcal C)$。
\begin{itemize}[leftmargin=1.6em,itemsep=2pt,topsep=2pt]
  \item $Tw(\mathcal C)$ 的对象是 $\mathcal C$ 的单侧扭转复形之集合。
  \item 对单侧扭转复形 $(\mathcal X,\Phi),(\mathcal Y,\Psi)\in Tw(\mathcal C)$，态射空间取为相应的 $A_\infty$ 范畴 $\tilde{\mathcal C}$ 者：
  \[
    \bigl(Tw(\mathcal C)\bigr)\bigl((\mathcal X,\Phi),(\mathcal Y,\Psi)\bigr):=\tilde{\mathcal C}(\mathcal X,\mathcal Y).
  \]
  \item $Tw(\mathcal C)$ 的 $A_\infty$ 结构 $m_{n}^{Tw}$（$n=1,2,\dots$）定义为
  \begin{align*}
    &m_{n}^{Tw}(\varphi_{a_{1}a_{2}},\dots,\varphi_{a_{n}a_{n+1}}):={}\\
    &\qquad\sum_{k_{1},\dots,k_{n+1}\in\Z_{\ge 0}}\tilde m_{n+k_{1}+\cdots+k_{n+1}}\Bigl((\Phi_{a_{1}})^{k_{1}},\varphi_{a_{1}a_{2}},(\Phi_{a_{2}})^{k_{2}},\dots,(\Phi_{a_{n}})^{k_{n}},\varphi_{a_{n}a_{n+1}},(\Phi_{a_{n+1}})^{k_{n+1}}\Bigr).
  \end{align*}
  其中 $a_{i}:=(\mathcal X_{a_{i}},\Phi_{a_{i}})\in Tw(\mathcal C)$，$\varphi_{a_{i}a_{i+1}}\in(sTw(\mathcal C))(a_{i},a_{i+1})$。
\end{itemize}
\end{definition}

\begin{remark}{7.48}
如此得到的 $A_\infty$ 范畴 $Tw(\mathcal C)$，若原来的 $A_\infty$ 范畴 $\mathcal C$ 的高阶积平凡（$m_{3}=m_{4}=\cdots=0$），就与 5.5.2 节中由 DG 范畴 $\mathcal C$ 构成的 DG 范畴 $Tw(\mathcal C)$ 一致。特别地，$Tw(\mathcal C)$ 的积 $m_{2}^{Tw}$ 完全不受来自 $\Phi_{a_{i}}$ 的修正，与 $\tilde{\mathcal C}$ 中的积 $\tilde m_{2}$ 一致，这一点可以看出。
\end{remark}

\begin{definition}{7.49}
（三角范畴 $Tr(\mathcal C)$）对 $A_\infty$ 范畴 $\mathcal C$，把三角范畴 $Tr(\mathcal C)$ 定义为 $Tw(\mathcal C)$ 的零次上同调（定义 7.38）：
\[
  Tr(\mathcal C):=H^{0}\bigl(Tw(\mathcal C)\bigr).
\]
\end{definition}

\begin{definition}{7.50}
对三角范畴 $\mathcal T$，若存在 $A_\infty$ 范畴 $\mathcal C$ 使 $\mathcal T=Tr(\mathcal C)$，则称 $A_\infty$ 范畴 $Tw(\mathcal C)$ 为三角 $A_\infty$ 范畴（triangulated $A_\infty$-category），并称 $\mathcal T$ 具有 $A_\infty$ 增广（$A_\infty$-enhancement）$Tw(\mathcal C)$。
\end{definition}

% [原文 p.188]
$\mathcal C$ 是 $Tw(\mathcal C)$ 的全子 $A_\infty$ 范畴，可以看作 $\mathcal T=Tr(\mathcal C)$ 的生成系。

$Tr(\mathcal C)$ 的正合三角结构如下确定。首先，对 $Tw(\mathcal C)$ 中的态射
\[
  \Psi\in\bigl(sTw(\mathcal C)\bigr)^{0}\bigl((\mathcal Z,\Phi_{\mathcal Z}),(\mathcal X,\Phi_{\mathcal X})\bigr),
\]
令
\[
  \mathcal Y:=\mathcal Z\oplus\mathcal X,\qquad
  \Phi_{\mathcal Y}:=\begin{pmatrix}\Phi_{\mathcal Z}&\Psi\\ 0&\Phi_{\mathcal X}\end{pmatrix},
\]
则 $(\mathcal Y,\Phi_{\mathcal Y})$ 是单侧扭转复形。实际上，它成为单侧扭转复形的条件与 $m_{1}^{Tw}(\Psi)=0$ 等价。由此，把 $A_\infty$ 范畴 $sTw(\mathcal C)$ 的次数 $-1$ 的态射所成的系列
\[
  \cdots\ (\mathcal Z,\Phi_{\mathcal Z})[-1]\xrightarrow{w[-1]}(\mathcal X,\Phi_{\mathcal X})\xrightarrow{u}(\mathcal Y,\Phi_{\mathcal Y})\xrightarrow{v}(\mathcal Z,\Phi_{\mathcal Z})\xrightarrow{w}\cdots
\]
中的 $w[-1]$ 定义为由 $\Psi\in(sTw(\mathcal C))^{0}((\mathcal Z,\Phi_{\mathcal Z}),(\mathcal X,\Phi_{\mathcal X}))$ 诱导的态射，$u,v$ 分别定义为由恒等态射 $\id_{\mathcal X},\id_{\mathcal Z}$ 自然确定的态射。把由这样的系列诱导出的 $Tr(\mathcal C)$ 中的（由次数零的态射构成的）系列同型的系列定为正合三角序列，则（与阿贝尔范畴中复形的同伦范畴等情形同样地）可以证明它们确定出正合三角结构。

% [Tutor Note: 这里的构造与 5.2 节 $K(\mathcal C)$ 的锥序列构造完全平行：$(\mathcal Y,\Phi_{\mathcal Y})$ 就是 $\Psi$ 的「锥」，$u,v$ 是自然投影与嵌入。要点在于 $\Psi$ 须满足 $m_1^{Tw}(\Psi)=0$（即闭链），否则 $\Phi_{\mathcal Y}$ 不满足 MC 方程。可验证：$\mathcal C$ 为 DG 范畴时常数项退化，得到 5.5.2 节的锥序列。]

\subsection{该构造下三角范畴中的推论}

\begin{theorem}{7.51}
若两个 $A_\infty$ 范畴 $\mathcal C,\mathcal C'$ 为 $A_\infty$ 同伦等价，则 $Tr(\mathcal C)\simeq Tr(\mathcal C')$。
\end{theorem}

在叙述这个证明之前，先叙述由该定理得到的若干事实。首先，由形式性的定义（定义 7.19、7.43）以下事实是显然的。

\begin{corollary}{7.52}
当 $A_\infty$ 范畴 $\mathcal C$ 是形式的时，
\[
  Tr(\mathcal C)\simeq Tr(H(\mathcal C))
\]
成立。
\end{corollary}

也就是说，$Tr(\mathcal C)$ 仅由 $\mathcal C$ 的上同调信息决定。作为更特殊的情形，得到以下结论。

\begin{corollary}{7.53}
若 $A_\infty$ 范畴 $\mathcal C$ 对所有对象 $X,Y\in\mathcal C$ 与整数 $r\ne 0$ 都满足 $H^{r}(\mathcal C)(X,Y)=0$，则
\[
  Tr(\mathcal C)\simeq Tr\bigl(H^{0}(\mathcal C)\bigr)
\]
成立。
\end{corollary}

% [Tutor Note: 系 7.53 按作者正誤表（2014-05-20 版）p.188 条订正：底本原印作 $Tr(\mathcal C)\simeq Tr(H(\mathcal C))$，应为 $Tr(\mathcal C)\simeq Tr(H^{0}(\mathcal C))$。理由：条件 $H^r(\mathcal C)=0$（$r\neq0$）意味着 $H(\mathcal C)$ 与 $H^0(\mathcal C)$ 作为分次范畴只差零次部分，故 $H(\mathcal C)$ 的增广退化为 $H^0(\mathcal C)$ 的增广。本译本已按订正后译出。]

特别地，当 $\mathcal C$ 的对象只有有限个时，$Tr(\mathcal C)$ 仅由代数
\[
  A:=\bigoplus_{X,Y\in\Ob(\mathcal C)}H^{0}\bigl(\mathcal C(X,Y)\bigr)
\]
的信息决定。

% [原文 p.189]
\begin{proof}
由态射的次数计算，可得 $A_\infty$ 范畴 $\mathcal C$ 是形式的。实际上，只需考虑 $\mathcal C$ 为极小 $A_\infty$ 范畴的情形即可；当 $k\ge 3$ 时，对 $X_{i_{1}},\dots,X_{i_{k+1}}\in\mathcal C$、$a_{j(j+1)}\in\mathcal C^{0}(X_{i_{j}},X_{i_{j+1}})$（$j=1,\dots,k$）有
\[
  \bigl|m_{k}(a_{12},\dots,a_{k(k+1)})\bigr|=2-k\ne 0,
\]
故非平凡的高阶 $A_\infty$ 积不可能存在。\hfill$\square$
\end{proof}

% [Tutor Note: 这段证明是「次数逼迫」的标准论证：$m_k$ 的次数为 $2-k$，而极小范畴中各 $\mathcal C^0(X,Y)$ 的态射次数为 $0$，故 $k\ge3$ 时 $m_k$ 的像次数为负，只能为零。补题：此时 $m_2$ 已满足结合律，故 $\mathcal C$ 与其上同调范畴同构，即形式。]

$\mathcal C$ 的对象构成 $Tr(\mathcal C)$ 的强例外对象系时，就是系 7.53 的特殊情形。由此可知：若具有 DG 增广的 $K$ 上三角范畴 $\mathcal T$ 具有强例外生成系 $(E_{1},\dots,E_{n})$，则 $\mathcal T$ 与自同态代数 $A:=\Hom_{\mathcal T}(E,E)$（$E:=\bigoplus_{i=1}^{n}E_{i}$）上的有限生成右模的导出范畴 $D^{b}(\mathrm{mod}\text{-}A)$ 三角范畴等价（定理 6.53）。

$A_\infty$ 范畴 $\mathcal C$ 的对象在 $Tr(\mathcal C)$ 中构成例外对象系时，称 $\mathcal C$ 为有序 $A_\infty$ 范畴（directed $A_\infty$-category）。此时所得的三角范畴 $Tr(\mathcal C)$ 含有大量结构处于可理解水平的例子。特别地，当其例外对象系是强例外对象系时，如上所述就归结为有限维代数的表示论。

\textbf{定理 7.51 的证明概略}（cf.\ [28]）：$A_\infty$ 范畴 $\mathcal C$ 中的单侧扭转复形 $(\mathcal X,\Phi)\in Tw(\mathcal C)$，是由对象 $\mathcal X=\bigoplus_{i=1}^{l}\tilde X_{i}$ 上 $A_\infty$ 范畴 $\tilde{\mathcal C}$ 的 $A_\infty$ 结构所对应的 $A_\infty$ Maurer--Cartan 方程的解 $\Phi$ 给出的。不过 $\Phi=\{\phi_{ij}\}$ 是对角成分为零的上三角矩阵（$i\ge j$ 时 $\phi_{ij}=0$）。由这个矩阵的形状，$\Phi$ 的分量具有以 $i-j=1,\dots,l-1$ 为次数的滤过结构。（可以想成变形理论中形式参数 $h$ 的 $i-j$ 次幂在起作用。）对每个 $\mathcal X$，都能定义其 $A_\infty$ Maurer--Cartan 方程的解集 $\mathrm{MC}_{\tilde{\mathcal C}}(\mathcal X)$，而三角范畴 $Tr(\mathcal C)$ 的对象，就是每确定一个 $\mathcal X$，与 $\mathrm{MC}_{\tilde{\mathcal C}}(\mathcal X)$ 的一个元素相对应。进而在集合 $\mathrm{MC}_{\tilde{\mathcal C}}(\mathcal X)$ 上可以定义规范等价关系 $\overset{\mathrm{g.e.}}{\sim}$，其商集记作
\[
  \mathcal M_{\tilde{\mathcal C}}(\mathcal X):=\mathrm{MC}_{\tilde{\mathcal C}}(\mathcal X)\big/\overset{\mathrm{g.e.}}{\sim}.
\]

在特征零的体 $K$ 上的 $A_\infty$ 范畴的情形，系 7.35 断言：对 $A_\infty$ 同伦等价的 $A_\infty$ 范畴 $\mathcal C,\mathcal C'$ 与每个对象 $\mathcal X\in\Ob(\tilde{\mathcal C})=\Ob(\tilde{\mathcal C}')$，集合 $\mathcal M_{\tilde{\mathcal C}}(\mathcal X)$ 与 $\mathcal M_{\tilde{\mathcal C}'}(\mathcal X)$ 之间存在双射。以此为基础，可以证明以下各点。
\begin{itemize}[leftmargin=1.5em,itemsep=2pt,topsep=2pt]
  \item 对 $A_\infty$ 范畴 $\mathcal C$，若 $\mathrm{MC}(\mathcal X)$ 的两个元素规范等价，则它们作为三角范畴 $Tr(\mathcal C)$ 的对象同构。
  \item $A_\infty$ 写像 $\mathbf f\colon\mathcal C\to\mathcal C'$ 诱导 $A_\infty$ 写像 $\mathbf f^{Tw}\colon Tw(\mathcal C)\to Tw(\mathcal C')$，由此得到的函子 $f^{Tr}\colon Tr(\mathcal C)\to Tr(\mathcal C')$ 是三角函子。
  \item 若 $\mathbf f,\mathbf g\colon\mathcal C\to\mathcal C'$ 为 $A_\infty$ 同伦，则 $f^{Tr}=g^{Tr}$。
\end{itemize}

% [原文 p.190]
把这些合起来，就得到对应于系 7.35 之推广的定理 7.51。

当体 $K$ 的特征不为零时，一般地「$A_\infty$ 函子是 $A_\infty$ 拟同构函子」与「它是 $A_\infty$ 同伦等价」并不等价；但可直接应用关于 $A_\infty$ 拟同构函子与 $A_\infty$ 范畴的分解定理（定理 7.21）来修正上述讨论，从而得到定理 7.51（细节见 [28]）。

\section{$A_\infty$ 模的范畴与米田嵌入}

对 $A_\infty$ 范畴 $\mathcal C$，可以考虑 $\mathcal C$ 上 $A_\infty$ 模所构成的 DG 范畴 $\mathrm{Mod}\text{-}\mathcal C$。此时存在称为米田嵌入的 $A_\infty$ 函子 $\mathbf f\colon\mathcal C\to\mathrm{Mod}\text{-}\mathcal C$。据此，对具有恒等元的 $A_\infty$ 范畴 $\mathcal C$，有构造与其 $A_\infty$ 拟同构的、具有恒等元的 DG 范畴的方法。本章对此进行说明。7.4.1 节讨论 $A_\infty$ 代数 $A$ 的米田嵌入，7.4.2 节把它们推广到 $A_\infty$ 范畴的情形。

\subsection{$A_\infty$ 代数的米田嵌入}

定义 $A_\infty$ 代数 $(A[1],\mathbf m)$ 上的模所构成的 DG 范畴 $\mathrm{Mod}\text{-}(A[1],\mathbf m)$。

\begin{definition}{7.54}
（$A_\infty$ 模（$A_\infty$-module））给定 $A_\infty$ 代数 $(A[1],\mathbf m)$ 时，右 $(A[1],\mathbf m)$ 模 $(M,m^{M})$ 是指：某个 $\Z$ 分次向量空间 $M$，与次数为 $1$ 的多重线性写像的族
\[
  \{m_{1+k}^{M}\colon M\otimes(A[1])^{\otimes k}\to M\}_{k\ge 0},
\]
且对 $n=1,2,\dots$ 满足以下关系式者：
\begin{equation}
\begin{aligned}
  0={}&\sum_{k+l=n}m_{1+k}^{M}\bigl(m_{1+l}^{M}(\xi,a_{1},\dots,a_{l}),a_{l+1},\dots,a_{n}\bigr)\\
  &+\sum_{k+l=n+1}(-1)^{|\xi|+|a_{1}|+\cdots+|a_{j}|}\,m_{1+k}^{M}\bigl(\xi,a_{1},\dots,a_{j},m_{l}(a_{j+1},\dots,a_{j+l}),a_{j+l+1},\dots,a_{n}\bigr).
\end{aligned}
\tag{7.8}
\end{equation}
\end{definition}

若 $(A,\mathbf m)$ 是（分次）代数（即 $m_{1}=m_{3}=m_{4}=\cdots=0$），并且对 $k=0,2,3,\dots$ 有 $m_{1+k}^{M}=0$，则（通过适当地调节次数）右 $(A[1],\mathbf m)$ 模 $M$ 不过是（分次）右 $A$ 模。此时特别地存在自由模 $M=A$。这有 $A_\infty$ 模的类似物。

\begin{proposition}{7.55}
令 $M=A[1]$，取 $m_{1+k}^{M}=m_{1+k}$（即 $\{m_{1+k}^{M}\}_{k\ge 0}$），则 $(M,m^{M})$ 构成右 $(A[1],\mathbf m)$ 模。
\end{proposition}

与 $A_\infty$ 代数的棒构成（7.1.2 节）同样，对右 $(A[1],\mathbf m)$ 模 $(M,m^{M})$，可以把 $\mathbf m$、$m^{M}$ 提升为 $M\otimes T^{c}(A[1])$ 上次数为 $1$ 的写像。把它记为 $\mathbf m^{M}+\mathbf m\colon M\otimes T^{c}(A[1])\to M\otimes T^{c}(A[1])$，则 $A_\infty$ 模的关系式

% [原文 p.191]
(7.8) 可以表为
\[
  (\mathbf m^{M}+\mathbf m)^{2}=0.
\]

\begin{definition}{7.56}
对两个右 $(A[1],\mathbf m)$ 模 $M,N$ 与整数 $k\ge 0$，把次数为 $r$ 的多重线性写像
\[
  M\otimes(A[1])^{\otimes k}\longrightarrow N
\]
的空间记作 $\Hom^{r,k}(M,N)$，并令 $\Hom^{r}(M,N):=\prod_{k=0}^{\infty}\Hom^{r,k}(M,N)$。
\end{definition}

元素 $f\in\Hom^{r,k}(M,N)$ 的自然提升 $\tilde f\colon M\otimes T^{c}(A[1])\to N\otimes T^{c}(A[1])$，通过：对 $\xi\in M$、$a_{1},\dots,a_{n}\in A[1]$ 令
\[
  \tilde f(\xi\otimes a_{1}\otimes\cdots\otimes a_{n}):=f(\xi,a_{1},\dots,a_{k})\otimes a_{k+1}\otimes\cdots\otimes a_{n}
\]
而确定。不过当 $k>n$ 时右边取为零。进而可知
\begin{equation}
  D(f):=-(\mathbf m^{N}+\mathbf m)\circ\tilde f+(-1)^{r}\tilde f\circ(\mathbf m^{M}+\mathbf m) \tag{7.9}
\end{equation}
确定了 $\Hom^{r+1}(M,N)$ 的元素之自然提升。把对应的元素记作 $D(f)\in\Hom^{r+1}(M,N)$。

\begin{proposition}{7.57}
由对应 $f\mapsto D(f)$ 所确定的多重线性写像 $D\colon\Hom^{r}(M,N)\to\Hom^{r+1}(M,N)$ 满足 $D^{2}=0$。
\end{proposition}

\begin{proof}
由 $(\mathbf m^{M}+\mathbf m)^{2}=0$、$(\mathbf m^{N}+\mathbf m)^{2}=0$ 即得。\hfill$\square$
\end{proof}

由此确定的复形记作 $(\Hom(M,N),D)$。

\begin{definition}{7.58}
如下定义右 $(A[1],\mathbf m)$ 模所构成的 DG 范畴 $\mathrm{Mod}\text{-}(A[1],\mathbf m)$。
\begin{itemize}[leftmargin=1.6em,itemsep=2pt,topsep=2pt]
  \item 对象取右 $(A[1],\mathbf m)$ 模 $M,N,\dots$。
  \item 对右 $(A[1],\mathbf m)$ 模 $M,N\in\mathrm{Mod}\text{-}(A[1],\mathbf m)$，态射空间取为
  \[
    \Hom_{\mathrm{Mod}\text{-}(A[1],\mathbf m)}(M,N):=\Hom(M,N).
  \]
  \item 态射空间的微分结构取为复形 $(\Hom(M,N),D)$ 的微分。
  \item 态射的积 $m\colon\Hom(N,E)\times\Hom(M,N)\to\Hom(M,E)$ 由对 $f\in\Hom(M,N)$、$g\in\Hom(N,E)$ 的合成 $m(g,f):=g\circ f$ 给出。
\end{itemize}
\end{definition}

这构成（按定义 5.49 的记法）DG 范畴是显然的；特别地，由这个 DG 结构，$H^{0}(\mathrm{Mod}\text{-}(A[1],\mathbf m))$ 构成三角范畴。

\begin{definition}{7.59}
把由补题 7.55 的对象 $A[1]$ 单独构成的 $\mathrm{Mod}\text{-}(A[1],\mathbf m)$ 的全 DG 子范畴记作 $\mathrm{End}(A[1])$。它实际上是 DG 代数。
\end{definition}

以下，把 DG 代数、DG 范畴作为 $A_\infty$ 代数、$A_\infty$ 范畴的特别情形来讨论。

\begin{theorem}{7.60}
给定 $A_\infty$ 代数 $(A[1],\mathbf m)$ 时，存在从 $(A,\mathbf m)$ 到 $\mathrm{End}(A[1])$ 的 $A_\infty$ 拟同构写像。
\end{theorem}

% [原文 p.192]
以下在本章中证明这个定理。首先，这个 $A_\infty$ 拟同构写像，通过考虑 $a_{1},\dots,a_{n},\xi,o_{1},\dots,o_{m}\in A[1]$ 上的 $A_\infty$ 关系式 (7.3) 来构成。把所得的等式分成 $0=(a+b+c)$ 三项。这里令
\begin{align*}
  a&:=\sum(-1)^{|a_{1}|+\cdots+|a_{j}|}\,m_{k}\bigl(a_{1},\dots,a_{j},m_{l}(a_{j+1},\dots,a_{j+l}),a_{j+l+1},\dots,a_{n},\xi,o_{1},\dots,o_{m}\bigr),\\
  b&:=\sum(-1)^{|a_{1}|+\cdots+|a_{j}|}\,m_{k}\bigl(a_{1},\dots,a_{j},m_{*}(a_{j+1},\dots,a_{n},\xi,o_{1},\dots,o_{l}),o_{l+1},\dots,o_{m}\bigr),\\
  c&:=\sum(-1)^{|a_{1}|+\cdots+|a_{n}|+|\xi|+|o_{1}|+\cdots+|o_{j}|}\,m_{k}\bigl(a_{1},\dots,a_{n},\xi,o_{1},\dots,o_{j},m_{l}(o_{j+1},\dots,o_{j+l}),o_{j+l+1},\dots,o_{m}\bigr).
\end{align*}
再把项 $b$ 分成 $b=b_{0}+b_{1}$ 两部分：
\begin{align*}
  b_{0}&:=\sum_{k+l=m}m_{1+k}\bigl(m_{n+1+l}(a_{1},\dots,a_{n},\xi,o_{1},\dots,o_{l}),o_{l+1},\dots,o_{m}\bigr)\\
  &\qquad+(-1)^{|a_{1}|+\cdots+|a_{n}|}\sum_{k+l}m_{k}\bigl(a_{1},\dots,a_{n},m_{1+l}(\xi,o_{1},\dots,o_{l}),o_{l+1},\dots,o_{m}\bigr),
\end{align*}
$b_{1}$ 为其余下的项。

现在如下定义保次数的多重线性写像 $f_{k}\colon(A[1])^{\otimes k}\to s(\Hom(M,M))$（$M=A[1]$）：
\begin{equation}
  f_{k}(a_{1},\dots,a_{k})=-s\Bigl(\sum_{l=0}^{\infty}m_{k+1+l}(a_{1},\dots,a_{k},*,*,\dots,*)\Bigr) \tag{7.10}
\end{equation}
其中 $s$ 是 DG 代数 $\mathrm{End}(A[1])$ 的悬垂。由此得
\begin{align}
  s\cdot a&:=-\sum_{k+l=n+1}(-1)^{|a_{1}|+\cdots+|a_{j}|}\,f_{k}\bigl(a_{1},\dots,a_{j},m_{l}(a_{j+1},\dots,a_{j+l}),a_{j+l+1},\dots,a_{n}\bigr)(\xi,o_{1},\dots,o_{m}), \tag{7.11}\\
  s\cdot b_{1}&:=\sum_{k+l=n}\bigl(m\bigl(f_{k}(a_{1},\dots,a_{k}),f_{l}(a_{k+1},\dots,a_{n})\bigr)\bigr)(\xi,o_{1},\dots,o_{m}), \tag{7.12}\\
  s\cdot(b_{0}+c)&=D\bigl(f_{n}(a_{1},\dots,a_{n})\bigr)(\xi,o_{1},\dots,o_{m}). \tag{7.13}
\end{align}
这里，(7.12) 中的 $m$ 是施以悬垂之后的东西，即
\[
  m(f,g):=(-1)^{|f|}s\bigl(s^{-1}(f)\circ s^{-1}(g)\bigr),
\]
而在 (7.13) 中也把 $sDs^{-1}$ 记作 $D$。

由此可知 (7.10) 的 $\{f_{k}\}_{k\ge1}$ 定义出 $A_\infty$ 写像 $\mathbf f\colon(A,\mathbf m)\to\mathrm{End}(A[1])$。因此，剩下只需看这个 $\mathbf f$ 是 $A_\infty$ 拟同构写像即可。为此，讨论 $A_\infty$ 写像 $\mathbf f$ 关于 $f_{1}$ 的性质。首先，微分


% [原文 p.193]
把 $D\colon\Hom^{r}(M,M)\to\Hom^{r+1}(M,M)$ 表为
\[
  D=D_{1}+D_{2}+\cdots,\qquad D_{l}\colon\Hom^{r,k}(M,M)\to\Hom^{r+1,k+l-1}(M,M).
\]
具体地，由 (7.9)，对 $\phi\in\Hom^{r,k}(M,M)$ 有
\begin{align*}
  (D_{1}\phi)(\xi,o_{1},\dots,o_{k})={}&-m_{1+0}^{M}\bigl(\phi(\xi,o_{1},\dots,o_{k})\bigr)+(-1)^{r}\phi\bigl(m_{1+0}^{M}(\xi),o_{1},\dots,o_{k}\bigr)\\
  &+\sum_{i=1}^{k}(-1)^{r+|\xi|+|o_{1}|+\cdots+|o_{i-1}|}\phi\bigl(\xi,o_{1},\dots,m_{1}(o_{i}),\dots,o_{k}\bigr),\\
  (D_{2}\phi)(\xi,o_{1},\dots,o_{k+1})={}&-m_{1+1}^{M}\bigl(\phi(\xi,o_{1},\dots,o_{k}),o_{k+1}\bigr)+(-1)^{r}\phi\bigl(m_{1+1}^{M}(\xi,o_{1}),o_{2},\dots,o_{k+1}\bigr)\\
  &+\sum_{i=1}^{k}(-1)^{r+|\xi|+|o_{1}|+\cdots+|o_{i-1}|}\phi\bigl(\xi,o_{1},\dots,m_{2}(o_{i},o_{i+1}),\dots,o_{k+1}\bigr)
\end{align*}
等等。由 $(D)^{2}=0$，得 $(D_{1})^{2}=0$、$D_{1}D_{2}+D_{2}D_{1}=0,\dots$。

以下的补题是剩余证明的关键。

\begin{proposition}{7.61}
定义写像 $\kappa\colon\Hom^{r,k+1}(M,M)\to\Hom^{r-1,k}(M,M)$（$M=A[1]$）如下：对 $\phi\in\Hom^{r,k+1}(M,M)$，若 $k>0$ 则
\[
  (\kappa\phi)(\xi,o_{1},\dots,o_{k}):=\phi(\id_{A[1]},\xi,o_{1},\dots,o_{k}),
\]
若 $k=0$ 则令 $\kappa\phi=0$。$\kappa$ 具有以下性质：
\begin{enumerate}[label=(\alph*),leftmargin=2em,itemsep=1pt,topsep=2pt]
  \item $D_{1}\kappa+\kappa D_{1}=0$，
  \item $(D_{2}\kappa+\kappa D_{2})\phi=\phi$\quad（$\phi\in\Hom^{r,k}(M,M)$，$k>0$）。\hfill$\square$
\end{enumerate}
\end{proposition}

\begin{proof}
(a)、(b) 都可由直接计算证明。\hfill$\square$
\end{proof}

\begin{proposition}{7.62}
关于 (7.10) 的 $A_\infty$ 写像 $\mathbf f:=\{f_{1},f_{2},\dots\}$，复形的写像 $f_{1}\colon(A[1],m_{1})\to(\Hom(M,M),D)$ 所诱导的上同调上的写像 $H(f_{1})\colon H(A[1])\to H(\Hom(M,M))$ 是单射。
\end{proposition}

\begin{proof}
对关于次数齐次的元素 $a\in(A[1])^{|a|}$，考虑 $f_{1}(a)\in\Hom^{|a|+1}(M,M)$。只需证明：若存在 $\phi\in\Hom^{|a|}(M,M)$ 使
\begin{equation}
  f_{1}(a)=D(\phi) \tag{7.14}
\end{equation}
则存在 $b\in(A[1])^{|a|-1}$ 使 $a=m_{1}(b)$。实际上，把 $\phi=\phi^{0}+\phi^{1}+\cdots$（$\phi^{i}\in\Hom^{|a|,i}(M,M)$）表示，并注意 (7.14) 中 $\Hom^{|a|+1,1}(M,M)$ 的部分，即得 $m_{2}(a,\ )=D_{1}\phi^{0}$。也就是说，对任意元素 $o\in A[1]$ 有
\[
  m_{2}(a,o)=(D_{1}\phi^{0})(o)=-m_{1}^{M}\bigl(\phi^{0}(o)\bigr)+(-1)^{|a|}\phi^{0}\bigl(m_{1}(o)\bigr).
\]
特别地考虑 $o=\id_{A}$ 的情形，得
\[
  (-1)^{|a|}a=-m_{1}^{M}\bigl(\phi^{1}(\id_{A})\bigr)=-m_{1}\bigl(\phi^{1}(\id_{A})\bigr),
\]
故令 $b=(-1)^{|a|+1}\phi^{1}(\id_{A})$ 即得 $a=m_{1}(b)$。\hfill$\square$
\end{proof}

% [原文 p.194]
\begin{proposition}{7.63}
写像 $H(f_{1})\colon H(A[1])\to H(\Hom(M,M))$ 是满射。
\end{proposition}

\begin{proof}
只需证明：当元素 $\phi\in\Hom^{r}(M,M)$ 满足 $D(\phi)=0$ 时，必存在 $a\in A[1]$ 与 $\psi\in\Hom^{r-1}(M,M)$ 使
\begin{equation}
  \phi-f_{1}(a)-D(\psi) \tag{7.15}
\end{equation}
为零。用归纳法证明。把 $\phi:=\phi^{0}+\phi^{1}+\cdots$（$\phi^{i}\in\Hom^{r,i}(M,M)$）表示，则由 $D(\phi)=0$ 得 $D_{1}(\phi^{0})=0$、$D_{2}(\phi^{0})+D_{1}(\phi^{1})=0,\dots$。特别地在 $(\id_{A[1]},o)$ 上考虑其中第二个等式，得
\begin{align*}
  0&=\bigl(D_{2}(\phi^{0})\bigr)(\id_{A[1]},o)+\bigl(D_{1}\phi^{1}\bigr)(\id_{A[1]},o)\\
  &=\Bigl(-m_{1+1}^{M}\bigl(\phi^{0}(\id_{A[1]}),o\bigr)+(-1)^{r}\phi^{0}\bigl(m_{1+1}^{M}(\id_{A[1]},o)\bigr)\Bigr)-\bigl(\kappa D_{1}\phi^{1}\bigr)(o)\\
  &=\Bigl(-m_{1+1}^{M}\bigl(\phi^{0}(\id_{A[1]}),o\bigr)-(-1)^{r}\phi^{0}(o)\Bigr)+\bigl(D_{1}\kappa\phi^{1}\bigr)(o),
\end{align*}
于是令
\begin{equation}
  a:=-(-1)^{r}\phi^{0}(\id_{A}) \tag{7.16}
\end{equation}
则 $\phi-f_{1}(a)+D(\kappa\phi^{1})$ 到 $\Hom^{r,0}(M,M)$ 的限制为零。

现设存在由 (7.16) 给出的 $a$ 与某个 $\psi$，使得对 $n=1,\dots,k$，(7.15) 到 $\Hom^{r,n}(M,M)$ 的限制为零。此时把 (7.15) 到 $\Hom^{r,n}(M,M)$ 的限制记作 $\varphi^{n}$，则 $D_{1}(\varphi^{k+1})=0$、$D_{2}(\varphi^{k+1})+D_{1}(\varphi^{k+2})=0,\dots$。令 $\varphi':=\varphi-D\bigl(\kappa(\varphi^{k+1}+\varphi^{k+2})\bigr)$，则其到 $\Hom^{r,n}(M,M)$ 的限制 $\varphi'^{n}$ 在 $n=1,\dots,k$ 时为零（注意 $D_{1}\kappa\varphi^{k+1}=-\kappa D_{1}\varphi^{k+1}=0$），并且
\begin{align*}
  \varphi'^{k+1}&=\varphi^{k+1}-D_{2}\kappa\varphi^{k+1}-D_{1}\kappa\varphi^{k+2}\\
  &=\varphi^{k+1}-(\id-\kappa D_{2})\varphi^{k+1}-D_{1}\kappa\varphi^{k+2}\\
  &=\kappa D_{2}\varphi^{k+1}-D_{1}\kappa\varphi^{k+2}\\
  &=\kappa\bigl(-D_{1}\varphi^{k+2}+D_{1}\varphi^{k+2}\bigr)=0.
\end{align*}
于是归纳法完成，证明了 (7.15) 成立。\hfill$\square$
\end{proof}

由补题 7.62、补题 7.63 可知 $f_{1}$ 是复形的拟同构写像，因此 $\mathbf f$ 为 $A_\infty$ 拟同构写像（定理 7.60）的证明至此完成。

\subsection{$A_\infty$ 范畴中的米田嵌入}

作为 $A_\infty$ 代数米田嵌入（定理 7.60）的自然推广，有 $A_\infty$ 范畴的米田嵌入。先把 $A_\infty$ 代数换成 $A_\infty$ 范畴时，前节所定义的各种东西的类似物逐一加以定义。

\begin{definition}{7.64}
设 $\mathcal C$ 为具有恒等元的 $A_\infty$ 范畴。若对每个对象 $X\in\mathcal C$ 给定 $\Z$ 分次向量空间 $M(X)$，并对对象 $X,X_{1},X_{2},\dots,X_{k}$（$k\ge 0$）给定次数为 $1$ 的多重线性写像

% [原文 p.195]
\[
  m_{1+k}^{M}\colon M(X)\otimes s\mathcal C(X,X_{1})\otimes\cdots\otimes s\mathcal C(X_{k-1},X_{k})\longrightarrow M(X_{k}),
\]
且它们满足右 $A_\infty$ 加群的定义式 (7.8)，则称 $(M,m^{M})$ 为 $\mathcal C$ 上的右加群。
\end{definition}

给定 $\mathcal C$ 上的右加群 $M,N$ 时，对每个 $X,X_{1},\dots,X_{k}\in\mathcal C$ 给出多重线性写像
\[
  M(X)\otimes s\mathcal C(X,X_{1})\otimes\cdots\otimes s\mathcal C(X_{k-1},X_{k})\longrightarrow N(X_{k})
\]
这种对应所成的分次向量空间记作 $\Hom(M,N)$。按多重线性写像的次数 $r$ 分解表为 $\Hom(M,N)=\bigoplus_{r\in\Z}\Hom^{r}(M,N)$。

与定义 7.56 同样，微分 $D\colon\Hom^{r}(M,N)\to\Hom^{r+1}(M,N)$ 有定义。具有这个微分结构的、由右 $\mathcal C$ 加群所构成的 DG 范畴记作 $\mathrm{Mod}\text{-}\mathcal C$。这里 $\Hom_{\mathrm{Mod}\text{-}\mathcal C}^{r}(M,N):=\Hom^{r}(M,N)$，积结构则由自然的合成确定。

\begin{proposition}{7.65}
设 $\mathcal C$ 为具有恒等元的 $A_\infty$ 范畴。对每个对象 $Y\in\mathcal C$，存在如下确定的右 $\mathcal C$ 加群 $M_{Y}$。
\begin{itemize}[leftmargin=1.6em,itemsep=1pt,topsep=1pt]
  \item 对对象 $X\in\mathcal C$，$M_{Y}(X):=s\mathcal C(Y,X)$。
  \item $A_\infty$ 加群结构 $m_{1+k}^{M_{Y}}$ 取为由 $s\mathcal C$ 的 $A_\infty$ 结构自然确定者：
  \[
    m_{1+k}^{M_{Y}}(\xi,a_{01},\dots,a_{(k-1)k})=m_{1+k}(\xi,a_{01},\dots,a_{(k-1)k}).\hfill\square
  \]
\end{itemize}
\end{proposition}

\begin{theorem}{7.66}
（$A_\infty$ 范畴的米田嵌入）存在如下给出的 $A_\infty$ 函子 $\mathbf f=\{f;f_{1},f_{2},\dots\}\colon\mathcal C\to\mathrm{Mod}\text{-}\mathcal C$。
\begin{itemize}[leftmargin=1.6em,itemsep=2pt,topsep=2pt]
  \item 对 $Y\in\Ob(\mathcal C)$，$f(Y):=M_{Y}$（补题 7.65）。
  \item 对 $X_{1},\dots,X_{k+1}\in\mathcal C$、$a_{i(i+1)}\in s\mathcal C(X_{i},X_{i+1})$，
  \[
    f_{k}(a_{12},\dots,a_{k(k+1)}):=\sum_{l=0}^{\infty}m_{k+1+l}(a_{12},\dots,a_{k(k+1)};\ ;,\dots,).
  \]
\end{itemize}
特别地，把由 $f(\Ob(\mathcal C))$ 的对象所构成的 $\mathrm{Mod}\text{-}\mathcal C$ 的全 DG 子范畴记作 $\mathrm{Im}(f)$，则 $\mathbf f$ 诱导出 $\mathcal C$ 与 $\mathrm{Im}(f)$ 之间的 $A_\infty$ 拟同构函子。
\end{theorem}

% [Tutor Note: 定理 7.66 的公式中，分号「;」表示「此处插入模元素」的空位（与式 (7.10) 的记号 $*$ 同义），这里 $m_{k+1+l}$ 的第 $k+1$ 个自变量的位置留给 $M(X)$ 的元素。$\mathrm{Im}(f)$ 的全 DG 子范畴中，对象是每个 $Y$ 对应的 $M_Y$，态射由 $m_{k+1+l}$ 的截断给出。]

\begin{remark}{7.67}
上述定理 7.66，是作为 $A_\infty$ 代数米田嵌入（定理 7.60）的自然推广而得到的；但即使 $A_\infty$ 范畴 $\mathcal C$ 由有限个对象构成，对 $A_\infty$ 代数 $(A[1],\mathbf m)=\bigoplus_{X,Y\in\mathcal C}s\mathcal C(X,Y)$ 应用定理 7.60 所得，也与上述定理 7.66 略有不同。也就是说，在定理 7.66 中 $\mathrm{Im}(f)$ 由 $\sharp(\Ob(\mathcal C))$ 个对象构成，而在定理 7.60 中则考虑对应于它们之直和的右 $(A[1],\mathbf m)$ 加群 $M=A[1]$ 上的 DG 代数（即由 $1$ 个对象构成的 DG 范畴）。
\end{remark}

作为定理 7.66 的推广，例如给定 $A_\infty$ 范畴 $\mathcal C$ 及其全子 $A_\infty$ 范畴 $\mathcal C'$ 时，可知可得 $A_\infty$ 写像 $\mathbf f\colon\mathcal C\to\mathrm{Mod}\text{-}\mathcal C'$。把它应用于 $\mathcal C\subset Tw(\mathcal C)$ 的情形，就得到有限维代数情形下定理 6.53 的推广。

\begin{corollary}{7.68}
设 $\mathcal C$ 为由有限个对象 $E_{1},\dots,E_{k}$ 构成的 $K$ 上 $A_\infty$ 范畴。此时，定义 $\mathbf f=\{f;f_{1},f_{2},\dots\}$ 如下：
\begin{itemize}[leftmargin=1.6em,itemsep=1pt,topsep=1pt]
  \item 对 $\mathcal Y\in Tw(\mathcal C)$，$f(\mathcal Y):=M_{\mathcal Y}$（补题 7.65）；
  \item 对 $a_{i(i+1)}\in s\mathcal C(\mathcal X_{i},\mathcal X_{i+1})$、$\mathcal X_{i}:=(\mathcal X_{i},\Phi_{\mathcal X_{i}})$，
  \[
    f_{k}(a_{12},\dots,a_{k(k+1)}):=\sum_{l=0}^{\infty}m_{k+1+l}^{Tw}(a_{12},\dots,a_{k(k+1)};\ ;,\dots,),
  \]
\end{itemize}
则它是一个 $A_\infty$ 函子 $\mathbf f\colon Tw(\mathcal C)\to\mathrm{Mod}\text{-}\mathcal C$，并满足以下性质：
\begin{itemize}[leftmargin=1.6em,itemsep=1pt,topsep=1pt]
  \item $\mathbf f$ 给出 $Tw(\mathcal C)$ 与由 $f(Tw(\mathcal C))$ 的对象所构成的全 DG 子范畴 $\mathrm{Im}(f)\subset\mathrm{Mod}\text{-}\mathcal C$ 之间的 $A_\infty$ 拟同构。
  \item $\mathbf f$ 确定从 $Tr(\mathcal C)$ 到 $H^{0}(\mathrm{Mod}\text{-}\mathcal C)$ 的三角函子，特别地给出三角范畴等价 $Tr(\mathcal C)\simeq H^{0}(\mathrm{Im}(f))$。\hfill$\square$
\end{itemize}
\end{corollary}

\begin{remark}{7.69}
一般地，即使 $A_\infty$ 范畴 $\mathcal C,\mathcal C'$ 为 $A_\infty$ 拟同构，右 $\mathcal C$ 加群的 DG 范畴 $\mathrm{Mod}\text{-}\mathcal C$ 与 $\mathrm{Mod}\text{-}\mathcal C'$ 也未必是 $A_\infty$ 拟同构的。
\end{remark}

\section{$A_\infty$ 代数与 $A_\infty$ 范畴的例子}

代数、DG 代数等曾是 $A_\infty$ 代数的特别例子；作为与它们不同的例子，举出高阶 $A_\infty$ 积 $m_{3},m_{4},\dots$ 非平凡的极小 $A_\infty$ 代数。

\begin{example}{7.70}
设 $A$ 为以 $e^{0}=\id_{A},e^{2},e^{5}$ 为基的分次向量空间，其中 $e^{i}\in A^{i}$。只要对这些基确定 $A_\infty$ 积，$A_\infty$ 结构就算定义了。令
\[
  m_{2}(e^{0},e^{i})=e^{i}=m_{2}(e^{i},e^{0}),\quad i=0,2,5,\qquad m_{3}(e^{2},e^{2},e^{2})=e^{5},
\]
并把其余 $m_{3}$ 以及 $m_{4},m_{5},\dots$ 取为平凡者，则构成极小 $A_\infty$ 代数。
\end{example}

看上面的例子就会明白：其余的例子自己可以要多少有多少地构成。

另一方面，对 DG 代数应用极小模型定理，也能得到极小 $A_\infty$ 代数。

\begin{example}{7.71}
设 $A$ 为以 $e^{0}=\id_{A},e^{2},e^{3},h^{4},dh^{4},e^{6}$ 为基的分次向量空间。这里 $e^{i}\in A^{i}$、$h^{4}\in A^{4}$、$dh^{4}\in A^{5}$。微分由 $d(e^{i})=0$、$d(dh^{4})=0$、$d(h^{4})=dh^{4}$ 定义。积 $m$ 定义为
\[
  m(e^{2},e^{3})=dh^{4},\qquad m(e^{2},h^{4})=e^{6},
\]
而含 $e^{0}$ 者按 $e^{0}$ 为恒等元的方式定义，其余取为平凡。这构成 DG 代数。

对它应用 7.1.4 节极小模型的具体构成法，得到以下的极小 $A_\infty$ 代数 $(A',\mathbf m')$。$A'$ 的基为 $e^{i}$（$i=0,2,3,6$），$e^{0}$ 是 $(A',\mathbf m')$ 的恒等元，不含 $e^{0}$ 的积 $m'_{2}$ 全部平凡，$m'_{3}(e^{2},e^{2},e^{3})=-e^{6}$，其余高阶 $A_\infty$ 积也全部平凡。非平凡的积 $m'_{3}(e^{2},e^{2},e^{3})=-e^{6}$ 对应于 DG 代数中称为 Massey 积的东西（cf.\ [28], [42]）。
\end{example}

% [Tutor Note: 例 7.71 是极小模型定理的具体演算：原始 DG 代数中 $m(e^2,e^3)=dh^4$，而 $dh^4$ 在极小模型里被「同伦掉」（不再是基元素），于是结合律的偏差被 $m'_3(e^2,e^2,e^3)=-e^6$ 吸收。可验证：$e^2\cdot(e^2\cdot e^3)=e^2\cdot dh^4$ 与 $(e^2\cdot e^2)\cdot e^3=0$ 之差由 $m'_3$ 给出，符号 $-$ 来自 Koszul 符号 $(-1)^{|e^2|}=1$ 与 $h$ 的次数。]

\begin{example}{7.72}
再考虑「对称性更高」的 DG 代数 $(A,d,m)$，如下

% [原文 p.197]
取之。$A$ 的基为 $e^{i}$（$i=0,2,3,6,7,9$）、$h^{4}$、$dh^{4}$，其中 $e^{i}\in A^{i}$、$h^{4}\in A^{4}$、$dh^{4}\in A^{5}$，微分由 $d(e^{i})=0$、$d(dh^{4})=0$、$d(h^{4})=dh^{4}$ 定义。把 $e^{0}$ 取为恒等元，不含 $e^{0}$ 的积 $m$ 取为
\begin{align*}
  &m(e^{2},e^{3})=m(e^{3},e^{2})=dh^{4},\qquad m(e^{2},h^{4})=m(h^{4},e^{2})=e^{6},\\
  &m(e^{i},e^{9-i})=m(h^{4},dh^{4})=m(dh^{4},h^{4})=e^{9},\qquad i=2,3,6,7,
\end{align*}
其余取为平凡。这样的 DG 代数，由于积可交换，称为 DG 交换代数。与它对应的极小 $A_\infty$ 代数也具有非平凡的积 $m'_{3}$（习题）。
\end{example}

像这样的 DG 交换代数的例子，在 D.~Sullivan 的有理同伦论（cf.\ [55] 及其参考文献）等处被大量考虑。

$A_\infty$ 代数是仅由 $1$ 个对象构成的 $A_\infty$ 范畴；由多个对象构成的极小 $A_\infty$ 范畴的例子也可以手算构造出许多。

\begin{example}{7.73}
考虑以 $X_{1},X_{2},X_{3},X_{4}$ 为对象的极小 $A_\infty$ 范畴 $\mathcal C$。对 $i=1,2,3,4$ 令 $\mathcal C(X_{i},X_{i}):=K\cdot\id_{X_{i}}$；对 $i=1,2,3$ 令 $\mathcal C(X_{i},X_{i+1}):=K\cdot\alpha_{i(i+1)}$，其中 $\alpha_{i(i+1)}$ 是次数 $1$ 的基元；对 $i=1,2$ 令 $\mathcal C(X_{i},X_{i+2})=0$，并令 $\mathcal C(X_{1},X_{4}):=K\cdot\beta_{14}$，其中 $\beta_{14}$ 是次数 $2$ 的基元。
\[
\begin{tikzcd}[row sep=0.6em, column sep=3.2em]
  \bullet_{1} \arrow[r,"\alpha_{12}"] & \bullet_{2} \arrow[r,"\alpha_{23}"] & \bullet_{3} \arrow[r,"\alpha_{34}"] & \bullet_{4}
\end{tikzcd}
\]
除含恒等态射的积以外令 $m_{2}$ 平凡，并令
\[
  m_{3}(\alpha_{12},\alpha_{23},\alpha_{34}):=\beta_{14},
\]
其余高阶 $A_\infty$ 积 $m_{3},m_{4},\dots$ 也取为平凡，则 $\mathcal C$ 构成极小 $A_\infty$ 范畴。
\end{example}

作为它的推广，也可以考虑由 $(n+1)$ 个对象构成、只有 $m_{n}$ 为非平凡积的极小 $A_\infty$ 范畴。

像这样不含非平凡 $A_\infty$ 关系式的例子容易构造。更复杂的例子请参见 [28], [29] 等。

其实例 7.5 的 $A_\infty$ 范畴 $\mathcal C$ 在有限维代数的表示论中也会出现。对箭图
\[
  \vec\Delta:=\Bigl(\bullet_{1}\ \underset{\alpha_{21}}{\overset{}{\rightleftarrows}}\ \bullet_{2}\ \underset{\alpha_{32}}{\overset{}{\rightleftarrows}}\ \bullet_{3}\ \underset{\alpha_{43}}{\overset{}{\rightleftarrows}}\ \bullet_{4}\Bigr)
\]
的关系式 $\rho:=(4\,|\,\alpha_{43}\alpha_{32}\alpha_{21}\,|\,1)$ 所带的道代数 $A:=(K\vec\Delta,\rho)$，三角范畴等价 $D^{b}(\mathrm{mod}\text{-}A)\simeq Tr(\mathcal C)$ 成立。实际上，$\mathrm{mod}\text{-}A\subset D^{b}(\mathrm{mod}\text{-}A)$ 的直不可分解投射对象构成 $D^{b}(\mathrm{mod}\text{-}A)$ 的强例外生成系 $(P(1),\dots,P(4))$，而单对象也构成例外生成系 $(S(4),\dots,S(1))$。$\mathcal C$ 是作为 $D^{b}(\mathrm{mod}\text{-}A)$ 的 $A_\infty$ 增广的、由 $X_{i}=S(i)$（$i=1,\dots,4$）构成的完全 $A_\infty$ 子范畴而得到的。结果，同态代数 $A=\bigoplus_{i,j=1}^{4}\Hom_{D^{b}(\mathrm{mod}\text{-}A)}(P(i),P(j))$ 中的关系式，在 $\mathcal C$ 中变成了高次 $A_\infty$ 积！（这样的对应可以理解为所谓 Koszul 对偶性的推广。）在这个例子中，直不可分解投射对象与单对象之间的关系，以及它们上面的 Serre 函子，都可以通过例外生成系的变异以相当不错的计算量具体地看出来。

% [Tutor Note: 「例 7.5」在底本 p.197 如此印出，但第 7 章的例编号从 7.70 起，此处应是指第 6 章中关于箭图 $\vec\Delta$ 与关系式 $\rho$ 的例（作者在 p.197 前文也用「実は例 7.5 の A∞ 圏 C」指代同一对象）。作者正誤表未收此条。按内容判断，它指 6.5 节中作为 $D^b(\mathrm{mod}\text{-}A)$ 的 $A_\infty$ 增广而出现的例子，本译本照印并指出。]

% [原文 p.198]
对怎样的 $A_\infty$ 范畴 $\mathcal C$，三角范畴 $Tr(\mathcal C)$ 具有怎样的结构？这类问题不用说是目前几乎未被研究的、有趣的问题；但如上所述，对于使某个代数 $A$ 满足三角范畴等价 $D^{b}(\mathrm{mod}\text{-}A)\simeq Tr(\mathcal C)$ 的那种 $A_\infty$ 范畴 $\mathcal C$，也有情形是 $\mathcal C$ 所确定的 $Tr(\mathcal C)$ 那边的三角范畴结构更容易理解，这一点今后可望有应用。

\medskip
\noindent\textbf{文献：}深谷 [13] 中含有关于 $A_\infty$ 代数的解说。关于 $A_\infty$ 代数的基本性质，见 [26] 及其参考文献等。Keller [35]、Seidel [49] 中有关于 $A_\infty$ 范畴的解说。7.1.5 节、7.1.6 节的内容基本上来自与寺嶋氏的共同研究 [31]。关于三角 $A_\infty$ 范畴中三角范畴的结构与 $A_\infty$ 结构之间的关系，[28]、[29] 中有更详细的讨论。那里也有 $A_\infty$ 范畴更复杂的例子。
